\chapter*{Anexos}
\phantomsection
\addcontentsline{toc}{chapter}{Anexos}

% Tablas y figuras de los anexos con numeración propia: A.1, A.2, etc.
\setcounter{table}{0}\renewcommand{\thetable}{A.\arabic{table}}
\setcounter{figure}{0}\renewcommand{\thefigure}{A.\arabic{figure}}

\section*{Participación en congresos}

{\setlength{\leftskip}{1cm}\noindent\fullcite{bajacfigueredo2024}\par}
{\setlength{\leftskip}{1cm}\noindent Trabajo presentado en el XLIII Congresso Nacional de Matemática Aplicada
e Computacional, Porto de Galinhas, Pernambuco, Brasil, del 16 al 20 de septiembre de 2024.\par}

% Este anexo empieza en página nueva: así la Tabla A.1 arranca junto a su texto y la nota no queda sola.
\clearpage
\section*{Contrastes de Wilcoxon completos}

La Tabla~\ref{tab:wilcoxon_completa} reúne los sesenta y seis contrastes de Wilcoxon de las seis métricas,
calculados sobre los veinte pares. En cada métrica hay once: seis comparan la propuesta con los demás
métodos y cinco comparan el Top-Hat clásico con cada método multiescala. La corrección de Holm se aplicó
dentro de cada métrica, sobre esos once contrastes. La columna Escenas repite la decisión con las trece
escenas como bloques, en lugar de los veinte pares. Los contrastes de la propuesta se resumen en la
Tabla~\ref{tab:wilcoxon_resumen}, de la Sección~\ref{sec:estadistico}.

{\setstretch{1}\setlength{\LTpost}{0pt}\setlength{\tabcolsep}{5pt}
\begin{longtable}{llcccc}
\caption{Contrastes de Wilcoxon pareados de las seis métricas, con corrección de Holm}\label{tab:wilcoxon_completa} \\
\toprule
& & & \textbf{\boldmath $p$} & \multicolumn{2}{c}{\textbf{Decisión}} \\
\cmidrule(lr){5-6}
\textbf{Método} & \textbf{Frente a} & \textbf{Diferencia} & \textbf{corregido} & \textbf{Pares} & \textbf{Escenas} \\
\midrule
\endfirsthead
\caption[]{Contrastes de Wilcoxon pareados de las seis métricas, con corrección de Holm (continuación)} \\
\toprule
& & & \textbf{\boldmath $p$} & \multicolumn{2}{c}{\textbf{Decisión}} \\
\cmidrule(lr){5-6}
\textbf{Método} & \textbf{Frente a} & \textbf{Diferencia} & \textbf{corregido} & \textbf{Pares} & \textbf{Escenas} \\
\midrule
\endhead
\bottomrule
\endfoot
\multicolumn{6}{l}{\textit{Entropía (EN)}} \\*
Propuesta & Top-Hat clásico & 0.0636 & $< 0.001$ & gana & gana \\*
Propuesta & Pirámide de Laplace & 0.1455 & 0.024 & gana & gana \\*
Propuesta & Pirámide de razón & 0.1758 & $< 0.001$ & gana & gana \\*
Propuesta & DWT (Haar) & 0.3040 & $< 0.001$ & gana & gana \\*
Propuesta & DTCWT & 0.2978 & $< 0.001$ & gana & gana \\*
Propuesta & DWT (db4) & 0.3410 & $< 0.001$ & gana & gana \\*
Top-Hat clásico & Pirámide de Laplace & 0.0820 & 0.189 & empata & empata \\*
Top-Hat clásico & Pirámide de razón & 0.1122 & 0.002 & gana & gana \\*
Top-Hat clásico & DWT (Haar) & 0.2404 & $< 0.001$ & gana & gana \\*
Top-Hat clásico & DTCWT & 0.2342 & $< 0.001$ & gana & gana \\*
Top-Hat clásico & DWT (db4) & 0.2775 & $< 0.001$ & gana & gana \\
\addlinespace \multicolumn{6}{l}{\textit{Desviación estándar (SD)}} \\*
Propuesta & Top-Hat clásico & 0.0087 & $< 0.001$ & gana & gana \\*
Propuesta & Pirámide de Laplace & $-$0.0111 & 0.231 & empata & empata \\*
Propuesta & Pirámide de razón & 0.0168 & $< 0.001$ & gana & gana \\*
Propuesta & DWT (Haar) & 0.0274 & $< 0.001$ & gana & gana \\*
Propuesta & DTCWT & 0.0267 & $< 0.001$ & gana & gana \\*
Propuesta & DWT (db4) & 0.0308 & $< 0.001$ & gana & gana \\*
Top-Hat clásico & Pirámide de Laplace & $-$0.0198 & 0.053 & empata & empata \\*
Top-Hat clásico & Pirámide de razón & 0.0081 & 0.005 & gana & empata \\*
Top-Hat clásico & DWT (Haar) & 0.0187 & $< 0.001$ & gana & gana \\*
Top-Hat clásico & DTCWT & 0.0180 & $< 0.001$ & gana & gana \\*
Top-Hat clásico & DWT (db4) & 0.0220 & $< 0.001$ & gana & gana \\
\addlinespace \multicolumn{6}{l}{\textit{Gradiente medio (MG)}} \\*
Propuesta & Top-Hat clásico & $-$0.0123 & $< 0.001$ & pierde & pierde \\*
Propuesta & Pirámide de Laplace & 0.0103 & $< 0.001$ & gana & gana \\*
Propuesta & Pirámide de razón & 0.0080 & $< 0.001$ & gana & gana \\*
Propuesta & DWT (Haar) & 0.0080 & $< 0.001$ & gana & gana \\*
Propuesta & DTCWT & 0.0100 & $< 0.001$ & gana & gana \\*
Propuesta & DWT (db4) & 0.0095 & $< 0.001$ & gana & gana \\*
Top-Hat clásico & Pirámide de Laplace & 0.0226 & $< 0.001$ & gana & gana \\*
Top-Hat clásico & Pirámide de razón & 0.0203 & $< 0.001$ & gana & gana \\*
Top-Hat clásico & DWT (Haar) & 0.0203 & $< 0.001$ & gana & gana \\*
Top-Hat clásico & DTCWT & 0.0223 & $< 0.001$ & gana & gana \\*
Top-Hat clásico & DWT (db4) & 0.0218 & $< 0.001$ & gana & gana \\
\addlinespace \multicolumn{6}{l}{\textit{Información mutua (MI)}} \\*
Propuesta & Top-Hat clásico & 0.1089 & $< 0.001$ & gana & gana \\*
Propuesta & Pirámide de Laplace & $-$0.6723 & $< 0.001$ & pierde & pierde \\*
Propuesta & Pirámide de razón & $-$0.0508 & 0.133 & empata & empata \\*
Propuesta & DWT (Haar) & $-$0.1224 & $< 0.001$ & pierde & pierde \\*
Propuesta & DTCWT & $-$0.1268 & $< 0.001$ & pierde & pierde \\*
Propuesta & DWT (db4) & $-$0.1342 & $< 0.001$ & pierde & pierde \\*
Top-Hat clásico & Pirámide de Laplace & $-$0.7812 & $< 0.001$ & pierde & pierde \\*
Top-Hat clásico & Pirámide de razón & $-$0.1597 & $< 0.001$ & pierde & pierde \\*
Top-Hat clásico & DWT (Haar) & $-$0.2313 & $< 0.001$ & pierde & pierde \\*
Top-Hat clásico & DTCWT & $-$0.2357 & $< 0.001$ & pierde & pierde \\*
Top-Hat clásico & DWT (db4) & $-$0.2431 & $< 0.001$ & pierde & pierde \\
\addlinespace \multicolumn{6}{l}{\textit{Similitud estructural (SSIM)}} \\*
Propuesta & Top-Hat clásico & 0.0944 & $< 0.001$ & gana & gana \\*
Propuesta & Pirámide de Laplace & $-$0.0474 & $< 0.001$ & pierde & pierde \\*
Propuesta & Pirámide de razón & $-$0.0467 & $< 0.001$ & pierde & pierde \\*
Propuesta & DWT (Haar) & $-$0.0429 & $< 0.001$ & pierde & pierde \\*
Propuesta & DTCWT & $-$0.0665 & $< 0.001$ & pierde & pierde \\*
Propuesta & DWT (db4) & $-$0.0578 & $< 0.001$ & pierde & pierde \\*
Top-Hat clásico & Pirámide de Laplace & $-$0.1418 & $< 0.001$ & pierde & pierde \\*
Top-Hat clásico & Pirámide de razón & $-$0.1411 & $< 0.001$ & pierde & pierde \\*
Top-Hat clásico & DWT (Haar) & $-$0.1373 & $< 0.001$ & pierde & pierde \\*
Top-Hat clásico & DTCWT & $-$0.1608 & $< 0.001$ & pierde & pierde \\*
Top-Hat clásico & DWT (db4) & $-$0.1522 & $< 0.001$ & pierde & pierde \\
\addlinespace \multicolumn{6}{l}{\textit{Relación señal-ruido de pico (PSNR)}} \\*
Propuesta & Top-Hat clásico & $-$0.0319 & 0.674 & empata & empata \\*
Propuesta & Pirámide de Laplace & 1.9009 & $< 0.001$ & gana & gana \\*
Propuesta & Pirámide de razón & $-$0.5279 & $< 0.001$ & pierde & pierde \\*
Propuesta & DWT (Haar) & $-$0.7477 & $< 0.001$ & pierde & pierde \\*
Propuesta & DTCWT & $-$0.7564 & $< 0.001$ & pierde & pierde \\*
Propuesta & DWT (db4) & $-$0.8114 & $< 0.001$ & pierde & pierde \\*
Top-Hat clásico & Pirámide de Laplace & 1.9328 & $< 0.001$ & gana & gana \\*
Top-Hat clásico & Pirámide de razón & $-$0.4960 & 0.001 & pierde & pierde \\*
Top-Hat clásico & DWT (Haar) & $-$0.7158 & $< 0.001$ & pierde & pierde \\*
Top-Hat clásico & DTCWT & $-$0.7245 & $< 0.001$ & pierde & pierde \\*
Top-Hat clásico & DWT (db4) & $-$0.7795 & $< 0.001$ & pierde & pierde \\
\end{longtable}
\nota{Elaboración propia. Cada métrica tiene once contrastes: la propuesta frente al Top-Hat clásico y cada uno de esos dos operadores frente a los cinco métodos multiescala. Diferencia: media del método de la columna Método menos media del método de la columna Frente a, sobre los veinte pares, en las unidades de cada métrica (decibelios en PSNR). $p$ corregido: valor $p$ con la corrección de Holm sobre los once contrastes de la métrica. Decisión: si el contraste es significativo al 5~\% después de la corrección, el método de la columna Método gana o pierde según el signo de la diferencia; si no lo es, empata. En las seis métricas, un valor mayor es mejor. Pares: decisión con los veinte pares del corpus como bloques, como en la Tabla~\ref{tab:wilcoxon_resumen}. Escenas: decisión con las trece escenas como bloques. Las métricas de cada escena se promedian sobre sus tomas, y se aplican la misma corrección y la misma regla. El Top-Hat clásico se usó con radio 5 y peso 1. Propuesta: el algoritmo de fusión presentado en este trabajo.}
}

